% Options for packages loaded elsewhere
% Options for packages loaded elsewhere
\PassOptionsToPackage{unicode}{hyperref}
\PassOptionsToPackage{hyphens}{url}
\PassOptionsToPackage{dvipsnames,svgnames,x11names}{xcolor}
%
\documentclass[
  letterpaper,
  DIV=11,
  numbers=noendperiod]{scrartcl}
\usepackage{xcolor}
\usepackage{amsmath,amssymb}
\setcounter{secnumdepth}{-\maxdimen} % remove section numbering
\usepackage{iftex}
\ifPDFTeX
  \usepackage[T1]{fontenc}
  \usepackage[utf8]{inputenc}
  \usepackage{textcomp} % provide euro and other symbols
\else % if luatex or xetex
  \usepackage{unicode-math} % this also loads fontspec
  \defaultfontfeatures{Scale=MatchLowercase}
  \defaultfontfeatures[\rmfamily]{Ligatures=TeX,Scale=1}
\fi
\usepackage{lmodern}
\ifPDFTeX\else
  % xetex/luatex font selection
\fi
% Use upquote if available, for straight quotes in verbatim environments
\IfFileExists{upquote.sty}{\usepackage{upquote}}{}
\IfFileExists{microtype.sty}{% use microtype if available
  \usepackage[]{microtype}
  \UseMicrotypeSet[protrusion]{basicmath} % disable protrusion for tt fonts
}{}
\makeatletter
\@ifundefined{KOMAClassName}{% if non-KOMA class
  \IfFileExists{parskip.sty}{%
    \usepackage{parskip}
  }{% else
    \setlength{\parindent}{0pt}
    \setlength{\parskip}{6pt plus 2pt minus 1pt}}
}{% if KOMA class
  \KOMAoptions{parskip=half}}
\makeatother
% Make \paragraph and \subparagraph free-standing
\makeatletter
\ifx\paragraph\undefined\else
  \let\oldparagraph\paragraph
  \renewcommand{\paragraph}{
    \@ifstar
      \xxxParagraphStar
      \xxxParagraphNoStar
  }
  \newcommand{\xxxParagraphStar}[1]{\oldparagraph*{#1}\mbox{}}
  \newcommand{\xxxParagraphNoStar}[1]{\oldparagraph{#1}\mbox{}}
\fi
\ifx\subparagraph\undefined\else
  \let\oldsubparagraph\subparagraph
  \renewcommand{\subparagraph}{
    \@ifstar
      \xxxSubParagraphStar
      \xxxSubParagraphNoStar
  }
  \newcommand{\xxxSubParagraphStar}[1]{\oldsubparagraph*{#1}\mbox{}}
  \newcommand{\xxxSubParagraphNoStar}[1]{\oldsubparagraph{#1}\mbox{}}
\fi
\makeatother


\usepackage{longtable,booktabs,array}
\usepackage{calc} % for calculating minipage widths
% Correct order of tables after \paragraph or \subparagraph
\usepackage{etoolbox}
\makeatletter
\patchcmd\longtable{\par}{\if@noskipsec\mbox{}\fi\par}{}{}
\makeatother
% Allow footnotes in longtable head/foot
\IfFileExists{footnotehyper.sty}{\usepackage{footnotehyper}}{\usepackage{footnote}}
\makesavenoteenv{longtable}
\usepackage{graphicx}
\makeatletter
\newsavebox\pandoc@box
\newcommand*\pandocbounded[1]{% scales image to fit in text height/width
  \sbox\pandoc@box{#1}%
  \Gscale@div\@tempa{\textheight}{\dimexpr\ht\pandoc@box+\dp\pandoc@box\relax}%
  \Gscale@div\@tempb{\linewidth}{\wd\pandoc@box}%
  \ifdim\@tempb\p@<\@tempa\p@\let\@tempa\@tempb\fi% select the smaller of both
  \ifdim\@tempa\p@<\p@\scalebox{\@tempa}{\usebox\pandoc@box}%
  \else\usebox{\pandoc@box}%
  \fi%
}
% Set default figure placement to htbp
\def\fps@figure{htbp}
\makeatother





\setlength{\emergencystretch}{3em} % prevent overfull lines

\providecommand{\tightlist}{%
  \setlength{\itemsep}{0pt}\setlength{\parskip}{0pt}}



 


\KOMAoption{captions}{tableheading}
\makeatletter
\@ifpackageloaded{caption}{}{\usepackage{caption}}
\AtBeginDocument{%
\ifdefined\contentsname
  \renewcommand*\contentsname{Table of contents}
\else
  \newcommand\contentsname{Table of contents}
\fi
\ifdefined\listfigurename
  \renewcommand*\listfigurename{List of Figures}
\else
  \newcommand\listfigurename{List of Figures}
\fi
\ifdefined\listtablename
  \renewcommand*\listtablename{List of Tables}
\else
  \newcommand\listtablename{List of Tables}
\fi
\ifdefined\figurename
  \renewcommand*\figurename{Figure}
\else
  \newcommand\figurename{Figure}
\fi
\ifdefined\tablename
  \renewcommand*\tablename{Table}
\else
  \newcommand\tablename{Table}
\fi
}
\@ifpackageloaded{float}{}{\usepackage{float}}
\floatstyle{ruled}
\@ifundefined{c@chapter}{\newfloat{codelisting}{h}{lop}}{\newfloat{codelisting}{h}{lop}[chapter]}
\floatname{codelisting}{Listing}
\newcommand*\listoflistings{\listof{codelisting}{List of Listings}}
\makeatother
\makeatletter
\makeatother
\makeatletter
\@ifpackageloaded{caption}{}{\usepackage{caption}}
\@ifpackageloaded{subcaption}{}{\usepackage{subcaption}}
\makeatother
\usepackage{bookmark}
\IfFileExists{xurl.sty}{\usepackage{xurl}}{} % add URL line breaks if available
\urlstyle{same}
\hypersetup{
  pdftitle={Oral GnRH antogonists for endometriosis-associated pain},
  colorlinks=true,
  linkcolor={blue},
  filecolor={Maroon},
  citecolor={Blue},
  urlcolor={Blue},
  pdfcreator={LaTeX via pandoc}}


\title{Oral GnRH antogonists for endometriosis-associated pain}
\usepackage{etoolbox}
\makeatletter
\providecommand{\subtitle}[1]{% add subtitle to \maketitle
  \apptocmd{\@title}{\par {\large #1 \par}}{}{}
}
\makeatother
\subtitle{A systematic review and meta-analysis}
\author{}
\date{}
\begin{document}
\maketitle


\section{Authors}\label{authors}

Mourad Elfaham1, Aya Attia1, Sara Ibrahim Abdelkader2*, Rahma Alaa
Abdelhafez2, Mohammed Ghanem2, Ahmed Azab2, Belgin Ahmed Nagah2, Nour
Atif2, Maha Moemen2, Ashraf Nabhan1,3

\textbf{Affiliation}

1 Faculty of Medicine, Ain Shams University, Cairo, Egypt.

2 Faculty of Medicine, Galala University, Attaka, Suez, Egypt.

3 Faculty of Medicine, MTI University, Cairo, Egypt.

*Corresponding author: Sara Ibrahim Abdelkader. Faculty of Medicine,
Galala University, Attaka, Suez, Egypt. Email: sara.elzeftawy@gu.edu.eg,
ORCID: 0009-0002-5877-4722

Word count: 2500

\section{Abstract}\label{abstract}

Background: Oral GnRH antagonists may improve endometriosis-related
pain. The benefits and harms of the available Oral GnRH antagonists
require a rigorous synthesis of evidence.

Objectives: To assess the effects of oral GnRH antagonists in women with
sendometriosis-related pain.

Search strategy: We searched CENTRAL, MEDLINE, Embase, and CINHAL on
January 6, 2025, using relevant search terms. Clinical trials registries
were searched for ongoing trials. The search had no language or date
restrictions.

Selection criteria: We included randomized controlled trials (RCTs)
comparing oral GnRH antagonists with placebo/no treatment/any
pharmacological intervention in women with sendometriosis-related pain.

Data collection and analysis: Two authors independently screened trials,
extracted data, and assessed the risk of bias in included studies using
Cochrane Risk of Bias tool 2. We used risk ratio (RR) for dichotomous
outcomes and mean difference (MD) for continuous outcomes, plus their
95\% confidence intervals (CIs). Pairwise meta-analyses were conducted,
where two or more studies met methodological criteria for inclusion.
GRADE approach was used to assess certainty of evidence for the
pre-specified outcomes.

Main results: We identified \#\# RCTs (\#\#\#\# participants). Five
studies (882 participants) compared oral GnRH antagonists with placebo.
OGA significantly achieved amenorrhea (RR 24.54; 95\% CI, 10.82-55.64),
reduced blood loss, and improved quality of life. The evidence for
adverse events is uncertain.

Conclusion: Compared with placebo, oral GnRH antagonists significantly
induces amenorrhea, reduces endometriosis associated pain, and improves
quality-of-life in women with endometriosis associated pain.

Systematic review registration: Center for Open Science, URL.

keywords: Gonadotropin-Releasing Hormone, Elagolix, Linzagolix,
Relugolix, Endometriosis, Dysmenorrhea, Quality of Life

\section{Background}\label{background}

Endometriosis, an enigmatic chronic health condition, affects 6-10\% of
women during their reproductive years. Women usually complains of
dysmenorrhea, chronic pelvic pain, deep dyspareunia and infertility.

The different types of pain associated with endometriosis can adversely
affect a woman's quality of life and her work productivity. Health
service utilization increases to cover diagnostic workup and the
provision of different lines of interventions.

Initial interventions to relief endometriosis-associated pain include
nonsteroidal anti-inflammatory drugs, depot medroxy progesterone
acetate, and Gonadotropin-releasing hormone (GnRH) agonists. Surgery for
women with symptomatic endometriosis offers an important line of
treatment at the time of a diagnostic laparoscopy or after an
insufficient response to medical therapy.

New short-acting nonpeptide oral GnRH antagonists recently emerged as an
alternative in the management of endometriois-associated pain. Oral GnRH
antagonists (Elagolix, Relugolix, and Linzagolix) act by suppressing the
anterior pituitary gonadotropin secretion. GnRH antagonists, as compared
to GnRH, do not cause an initial stimulation of the anterior pituitary
gland. Thus, this class does not result in an initial increase in
symptoms. The oral route would provide an additional advantage over GnRH
agonists. Oral GnRH antagonist produce a dose-dependent suppression of
ovarian estrogen production. The potential for Oral GnRH antagonists to
produce partial suppression at lower doses may decrease
endometriosis-related pain while avoiding the undesirable hypoestrogenic
effects like hot flashes, vaginal dryness, and the decrease in bone
mineral density.

The objective of this study was to assess the benefits and harms of Oral
GnRH antagonists in women with endometriosis-associated pain.

\section{Methods}\label{methods}

Protocol and registration

This systematic review was conducted following the methodological
standards of Cochrane Handbook {[}9{]}. We prospectively registered the
protocol in Open Science Platform (URL). The full text of the protocol
is available in an open access registry. We reported the review using
the Preferred Reporting Items for Systematic reviews and Meta-Analyses
(PRISMA) standards {[}10{]}. The full checklist is available as
Supplementary File 2.

Eligibility criteria

We included published randomized controlled trials that compared any
oral GnRH antagonist with no intervention, placebo or any active
pharmacological intervention in women with endometriosis associated pain
during the reproductive age, We assesse the followig outcomes:
non-menstrual pelvic pain, dysmenorrhea, dyspareunia, QOL, additional
use of analgesics, and adverse events (Any adverse event leading to
discontinuation, Bone mineral density, Hot flashes, Vaginal dryness).
The length of follow-up was up to 12 and 24 weeks.

Information sources

A comprehensive literature search was initially conducted on September
20, 2023, by two authors (AFN, HBA) who are information experts. We did
not impose language or other restrictions on any of the searches. We
searched bibliographic databases (Cochrane Central Register of
Controlled Trials (CENTRAL), MEDLINE/PubMed) and citation indexes (Web
of Science and Scopus). We included the terms (HELLP Syndrome) AND
(corticosteroids or glucocorticoids or Dexamethasone or Betamethasone or
Prednisolone). The search strategy was updated on February 3, 2024, and
new studies were not identified. The detailed exact strategy adapted for
each database is provided in Supplementary File 3 and is available as an
open access registry document. We peer-reviewed the search strategy
using PRESS checklist {[}11{]}, and further tested it with a set of
known relevant, `gold standard' reports. We also searched clinical trial
registries (ClinicalTrials.gov and the World Health Organization
International Clinical Trials Registry Platform) to identify ongoing
trials. We finally searched reference lists and explored the cited-by
logs of identified studies and previously published reviews.

Study selection

All reports identified in the databases were imported to Bibtex library
using Jabref version 5. After removing duplicates, two authors
independently screened all titles and abstracts for eligibility. We
retrieved and assessed the full text of all reports that potentially met
our eligibility criteria during screening. Two authors (AFK, HBA, MAE,
RHA) independently assessed each full-text article. Disagreements
regarding trial eligibility were resolved by consensus and finally
resolved by a third author (AFN).

Data collection process

For eligible studies, we extracted the data in duplicates using an
offline electronic form. We resolved discrepancies through discussion.
Extracted data were transcribed to a spreadsheet and checked for
accuracy. We contacted authors of the original reports, if needed, to
provide details regarding unclear or missing data.

Data items

Extracted data included study design, sample size, description of
included participants, description of the intervention, outcomes, trial
registration, and funding sources, and country.

Study risk of bias assessment

We assessed the following risk of bias domains as outlined in Cochrane
Handbook for Systematic Reviews of Interventions: (1) risk of bias
arising from the randomization process; (2) risk of bias due to
deviations from the intended interventions (effect of assignment to
intervention); (3) risk of bias due to missing outcome data; (4) risk of
bias in measurement of the outcome; and (5) risk of bias in selection of
the reported result. Each domain was judged as being at ``low risk of
bias'', ``some concerns'', or ``high risk of bias''. Trials with ``low
risk of bias'' in all domains were classified as being at overall ``low
risk of bias''. RCTs with one domain judged to be at ``some concerns'',
but no domain judged to be at ``high risk of bias'', were classified as
being at overall ``some concerns'' of risk of bias. RCTs were classified
as being at overall ``high risk of bias'' if at least one domain was
judged as being at ``high risk of bias''. However, if a trial was judged
to be at ``some concerns'' due to risk of bias for multiple domains, it
was judged as being at overall ``high risk of bias'' if the assessors
judged that the multiple concerns amounted to a serious risk of bias. In
case of discrepancies among their judgments and inability to reach
consensus, we consulted the senior author (AFN) to reach a final
decision.

Effect measures

For dichotomous data, we presented results as risk ratio (RR) with 95\%
confidence intervals (CI). For continuous data, we presented results as
mean diference (MD) with 95\% CI. The unit of analysis was the
individual participant. Data related to participants reported as not
compliant was analyzed on an intention-to-treat basis.

Synthesis methods

Common effect meta-analysis was performed to combine data of trials that
are judged to be sufficiently similar in terms of intervention,
populations, and methods. We planned to investigate substantial
statistical heterogeneity, defined as I² statistic ≥ 50\% or
P \textless{} 0.1.

We performed the planned subgroup analysis by type of antagonist
(Elagolix, Relugolix, Linzagolix). We assessed subgroup differences by
interaction tests. Results of the subgroup analyses were reported by
mentioning the Chi² statistic and P value, and the interaction test I²
value.

Sensitivity analysis was performed to explore robustness of pooled
estimate using outcome data from trials with a low risk of bias.

Synthesis was performed using RStudio 2024.12.0 Build 467 (MacOS, Apple
Silicon version), R 4.4 (2024-12-16) {[}12{]} and R package meta version
8 {[}13{]}.

Reporting bias assessment

We explored whether the study was included in a trial registry and
whether a protocol was available. We planned to examine funnel plots to
assess the potential for publication bias if we found 10 or more studies
reporting on a particular outcome.

Certainty assessment

We used the Grading of Recommendations, Assessment, Development and
Evaluation (GRADE) approach to create the Summary of Findings Table
{[}14{]} Briefly, GRADE uses study limitations, consistency of effect,
imprecision, indirectness, and publication bias to assess the certainty
of evidence for each outcome. A summary of the intervention effect and a
measure of certainty was produced using the GRADE Profiler Guideline
Development Tool (GRADEpro GDT) software {[}15{]} for the prespecified
important outcomes: maternal death, pulmonary edema, renal failure,
dialysis, liver morbidity, need for platelet transfusion, and perinatal
death. One author (AFN) conducted GRADE assessments and the decisions on
downgrading. This was discussed for final approval by all authors.

\section{Results}\label{results}

\section{Discussion}\label{discussion}

\section{Conclusions}\label{conclusions}

\section{List of abbreviations}\label{list-of-abbreviations}

If abbreviations are used in the text they should be defined in the text
at first use, and a list of abbreviations should be provided.

\section{Declarations}\label{declarations}

\subsection{Ethics approval and consent to
participate}\label{ethics-approval-and-consent-to-participate}

Not applicable.

\subsection{Consent for publication}\label{consent-for-publication}

Not applicable.

\subsection{Availability of data and
materials}\label{availability-of-data-and-materials}

All data relevant to this study are publicly available. Data, analysis
script and materials related to this study are publicly available on the
Open Science Framework at~{[}URL{]}. To facilitate reproducibility, this
manuscript was written by interleaving regular prose and analysis code
using R Markdown.

\subsection{Competing interests}\label{competing-interests}

The authors declare no competing interests.

\subsection{Funding}\label{funding}

This research received no specific grant from any funding agency in the
public, commercial or not-for-profit sectors.

\subsection{Authors' contributions}\label{authors-contributions}

AFN conceived the idea for this review and designed the review methods.
AFK, HBA coordinated the research activities. MAE, RRA, RHA, MNS, KMH,
RGE, RHH, MSE, MMA, YME, MYA, MHS contributed to search, screening,
interpretation of results. All authors collaborated in searching,
screening, selecting studies, data extraction and synthesis. AFN, AFK,
HBA wrote the first draft of the manuscript. All authors revised the
manuscript critically for important intellectual content. All authors
approved the final version of the manuscript. AFN is the guarantor for
this manuscript.

\subsection{Acknowledgements}\label{acknowledgements}

\section{References}\label{references}




\end{document}
